% from ega1-10-fr.tex
\oldpage[I]{190}

Soit $\mathscr{O}_{\mathfrak{X}}$ le faisceau d'anneaux topologiques
limite projective des faisceaux d'anneaux pseudo-discrets
$\mathscr{O}_i$, et soit
$u_i:\mathscr{O}_{\mathfrak{X}}\to\mathscr{O}_i$ l'homomorphisme
canonique. Alors l'espace topologiquement annelé
$(\mathfrak{X},\mathscr{O}_{\mathfrak{X}})$ est un préschéma formel~;
les homomorphismes $u_i$ sont surjectifs~; leurs noyaux
$\mathscr{J}^{(i)}$ forment un système fondamental d'Idéaux de définition
de $\mathfrak{X}$, et $\mathscr{J}^{(0)}$ est la limite projective des
faisceaux d'idéaux $\mathscr{J}_i$.
\end{proposition}

\begin{proof}
Notons d'abord que sur chaque fibre, $u_{ji}$ est un homomorphisme
surjectif et \emph{a fortiori} un homomorphisme local~; donc
$v_{ij}=(1_{\mathfrak{X}},u_{ji})$ est un morphisme de préschémas
$X_j\to X_i$ ($i\geq j$) (\hyperref[I.2.2.1-fr]{2.2.1}). Supposons
d'abord que chaque $X_i$ soit un schéma affine d'anneau $A_i$. Il existe
un homomorphisme \emph{d'anneaux}
$\varphi_{ji}:A_i\to A_j$ tel que
$u_{ji}=\widetilde{\varphi_{ji}}$
(\hyperref[I.1.7.3-fr]{1.7.3})~; par suite
(\hyperref[I.1.6.3-fr]{1.6.3}), le faisceau $\mathscr{O}_j$ est un
$\mathscr{O}_i$-Module quasi-cohérent sur $X_i$ (pour la loi externe
définie par $u_{ji}$), associé à $A_j$ considéré comme $A_i$-module à
l'aide de $\varphi_{ji}$. Pour tout $f\in A_i$, soit
$f'=\varphi_{ji}(f)$~; par hypothèse, les ouverts $D(f)$ et $D(f')$ sont
identiques dans $\mathfrak{X}$, et l'homomorphisme de
$\Gamma(D(f),\mathscr{O}_i)=(A_i)_f$ dans
$\Gamma(D(f),\mathscr{O}_j)=(A_j)_{f'}$ correspondant à $u_{ji}$ n'est
autre que $(\varphi_{ji})_f$ (\hyperref[I.1.6.1-fr]{1.6.1}). Mais
lorsqu'on considère $A_j$ comme $A_i$-module, $(A_j)_{f'}$ est le
$(A_i)_f$-module $(A_j)_f$, donc on a aussi
$u_{ji}=\widetilde{\varphi_{ji}}$ lorsque $\varphi_{ji}$ est cette fois
considéré comme homomorphisme de $A_i$-modules. Alors, comme $u_{ji}$ est
surjectif, on en conclut que $\varphi_{ji}$ l'est aussi
(\hyperref[I.1.3.9-fr]{1.3.9}) et si $\mathfrak{J}_{ji}$ est le noyau de
$\varphi_{ji}$, le noyau de $u_{ji}$ est un $\mathscr{O}_i$-Module
quasi-cohérent égal à $\widetilde{\mathfrak{J}_{ji}}$. En particulier,
on a $\mathscr{J}_i=\widetilde{\mathfrak{J}_i}$, où $\mathfrak{J}_i$ est
le noyau de $\varphi_{0i}:A_i\to A_0$. L'hypothèse b) entraîne que
$\mathscr{J}_i$ est \emph{nilpotent}~: en effet, comme $\mathfrak{X}$
est quasi-compact, on peut recouvrir $\mathfrak{X}$ par un nombre fini
d'ouverts $U_k$ tels que $(\mathscr{J}_i|U_k)^{n_k}=0$ et en prenant pour
$n$ le plus grand des $n_k$, on a $\mathscr{J}_i^n=0$. On en conclut que
$\mathfrak{J}_i$ est nilpotent (\hyperref[I.1.3.13-fr]{1.3.13}). Alors
l'anneau $A=\varprojlim A_i$ est admissible
(\hyperref[0.7.2.2-fr]{0, 7.2.2}), l'homomorphisme canonique
$\varphi_i:A\to A_i$ est surjectif et son noyau
$\mathfrak{J}^{(i)}$ est égal à la limite projective des
$\mathfrak{J}_{ik}$ pour $k\geq i$~; les $\mathfrak{J}^{(i)}$ forment
un système fondamental de voisinages de 0 dans $A$. Les assertions de
(\hyperref[I.10.6.3-fr]{10.6.3}) résultent dans ce cas de
(\hyperref[I.10.1.1-fr]{10.1.1}) et
(\hyperref[I.10.3.2-fr]{10.3.2}),
$(\mathfrak{X},\mathscr{O}_{\mathfrak{X}})$ n'étant autre que
$\operatorname{Spf}(A)$.

Toujours dans ce même cas particulier, notons que si $f=(f_i)$ est un
élément de la limite projective $A=\varprojlim A_i$, tous les ouverts
$D(f_i)$ (ouvert affine dans $X_i$) s'identifient à l'ouvert
$\mathfrak{D}(f)$ de $\mathfrak{X}$, le préschéma induit par $X_i$ sur
$\mathfrak{D}(f)$ s'identifiant donc au schéma affine
$\operatorname{Spec}((A_i)_{f_i})$.

Dans le cas général, remarquons d'abord que pour tout ouvert quasi-compact
$U$ de $\mathfrak{X}$, chacun des $\mathscr{J}_i|U$ est nilpotent, comme
le montre le raisonnement fait ci-dessus. Nous allons voir que pour tout
$x\in\mathfrak{X}$, il y a un voisinage ouvert $U$ de $x$ dans
$\mathfrak{X}$ qui est un \emph{ouvert affine pour tous les $X_i$}. En
effet, prenons $U$ ouvert affine pour $X_0$, et observons que
$\mathscr{O}_{X_0}=\mathscr{O}_{X_i}/\mathscr{J}_i$. Comme
$\mathscr{J}_i|U$ est nilpotent, en vertu de ce qui précède, $U$ est
ouvert affine aussi pour chaque $X_i$ en vertu de
(\hyperref[I.5.1.9-fr]{5.1.9}). Cela étant, pour tout $U$ satisfaisant
aux conditions précédentes, l'étude du cas affine faite plus haut montre
que $(U,\mathscr{O}_{\mathfrak{X}}|U)$ est un préschéma formel dont les
$\mathscr{J}^{(i)}|U$ forment un système fondamental d'idéaux de
définition et $\mathscr{J}^{(0)}|U$ est limite projective des
$\mathscr{J}_i|U$~; d'où la conclusion.
\end{proof}

\begin{corollary}[10.6.4]
\label{I.10.6.4-fr}
Supposons que pour $i\geq j$, le noyau de $u_{ji}$ soit
$\mathscr{J}_i^{j+1}$ et que $\mathscr{J}_1/\mathscr{J}_1^2$

